\documentclass[11pt,a4paper]{article}
\input{tutorial_style}

\begin{document}
\tutorialtitle{Implementation Companion}{Python Examples for Note 2 PINN/PIDL Cyber Control}{From loss terms to reproducible experiments}{Companion to the PINN/PIDL cyber-control tutorial note}
\tableofcontents
\newpage
\lhead{\small Note 2 Implementation Companion}

\section{Purpose}
This companion maps the Note 2 equations to the Python files in this repository. The aim is transparency: each example should make clear which unknown function is represented by a neural network, which residual is minimized, what data are needed, and what outputs should be inspected.

\section{File Map}
\begin{longtable}{p{0.34\textwidth}p{0.56\textwidth}}
\toprule
Python file & What it demonstrates \\
\midrule
\small\texttt{inverse\_pinn\_sir\_malware.py} & Inverse PINN: learn hidden S/I/R trajectories and unknown SIR parameters from sparse infected-state observations. \\
\small\texttt{pidl\_unknown\_mechanism.py} & PIDL correction: keep the known SIR mechanism explicit and learn a missing nonlinear transfer term. \\
\small\texttt{control\_pinn\_malware.py} & Direct neural-control PINN: train state and control networks using objective, residual, and initial-condition losses. \\
\small\texttt{pmp\_informed\_pinn\_malware.py} & PMP-informed PINN: train state, costate, and control networks with state, costate, stationarity, and boundary residuals. \\
\small\texttt{run\_training\_iterations.py} & Rebuilds CSV histories, baseline metrics, \texttt{OUTPUT\_PREVIEW.md}, \texttt{training\_summary.md}, and the diagnostic figures. \\
\bottomrule
\end{longtable}

\section{Common Implementation Pattern}
All four examples follow the same practical pattern:
\begin{enumerate}
\item define the cyber propagation model and parameter ranges;
\item choose the neural objects, such as state, correction, control, or costate networks;
\item sample observed data and collocation points;
\item build data, residual, boundary, objective, or PMP losses;
\item train with Adam and log each important loss component separately;
\item validate by inspecting trajectories, residuals, parameter estimates, and baseline comparisons.
\end{enumerate}

\section{Inverse PINN}
The inverse PINN observes sparse samples of the infected component \(I(t)\). The state network \(N_\theta(t)\) outputs \(\hat{x}(t)=[\hat{S},\hat{I},\hat{R}]\), and trainable positive parameters represent \(\hat{\beta}\) and \(\hat{\gamma}\).

\begin{methodbox}
\textbf{Loss structure.}
\[
\Loss = \Loss_{\rm data} + w_{\rm ic}\Loss_{\rm ic} + w_{\rm ode}\Loss_{\rm ode}.
\]
The data term matches sparse \(I(t)\) observations, the initial-condition term anchors \(x(0)\), and the ODE residual enforces SIR dynamics at collocation points.
\end{methodbox}

\begin{lstlisting}[style=pythonstyle,caption={Core inverse-PINN residual idea.}]
x_f = model(t_f)
dxdt = time_derivative(x_f, t_f)
rhs = sir_rhs(x_f, beta, gamma)
loss_ode = torch.mean((dxdt - rhs) ** 2)
loss = loss_data + w_ic * loss_ic + w_ode * loss_ode
\end{lstlisting}

\section{PIDL With A Missing Mechanism}
PIDL should not learn the entire dynamics when a trusted mechanism is already known. In this example the known SIR right-hand side remains explicit, and a correction network \(g_\omega(t,x)\) learns the missing nonlinear transfer.

\begin{equation}
\dot{x}=f_{\rm known}(x;\beta,\gamma)+B g_\omega(t,x).
\end{equation}

\begin{methodbox}
\textbf{What to inspect.}
The correction term should reduce residual error, but it should not become an unconstrained replacement for the known model. Inspect residual loss, correction regularization, and mean correction together.
\end{methodbox}

\section{Direct Control PINN}
The direct control PINN trains a state network \(\hat{x}_\theta(t)\) and a bounded control network \(\hat{u}_\phi(t)\). It minimizes a malware-control objective while enforcing the controlled SIR dynamics.

\begin{equation}
\Loss = J(\hat{x},\hat{u}) + w_{\rm res}\|\dot{\hat{x}}-f(\hat{x},\hat{u})\|^2 + w_{\rm ic}\|\hat{x}(0)-x_0\|^2.
\end{equation}

This method is simple and useful, but it is a direct optimization method. If a project claims PMP consistency, use the PMP-informed residuals below.

\section{PMP-Informed PINN}
The PMP-informed PINN adds a costate network \(\hat{\lambda}_\psi(t)\) and enforces the optimality system rather than only the state equation.

\begin{equation}
\Loss=w_x\|\dot{x}-f(x,u)\|^2
    +w_\lambda\|\dot{\lambda}+H_x(x,u,\lambda)\|^2
    +w_u\|H_u(x,u,\lambda)\|^2
    +w_b\Loss_{\rm boundary}.
\end{equation}

\begin{methodbox}
\textbf{Why the stationarity residual matters.}
A small state residual says the trajectory satisfies the dynamics. A small stationarity residual gives stronger evidence that the learned control is compatible with PMP necessary conditions.
\end{methodbox}

\section{Training Diagnostics}
The diagnostic figure generated by \texttt{scripts/run\_training\_iterations.py} has four panels:
\begin{itemize}
\item sparse-data inverse PINN: total loss and ODE residual;
\item PIDL missing mechanism: residual loss and learned correction size;
\item direct control PINN: objective, state residual, and mean control;
\item PMP-informed PINN: state, costate, stationarity, and total residuals.
\end{itemize}
Start with \texttt{artifacts/experiments/OUTPUT\_PREVIEW.md}, then open the CSV file behind the panel that matters for the method claim.

\section{Baseline Comparisons}
Training curves answer whether a loss is decreasing. Baseline comparisons answer whether the trained method is better than a simple alternative on a meaningful metric. The script therefore also writes \texttt{artifacts/experiments/baseline\_comparison\_metrics.csv} and \texttt{artifacts/figures/baseline\_comparison.png}.

\begin{longtable}{p{0.30\textwidth}p{0.32\textwidth}p{0.28\textwidth}}
\toprule
Topic & Baseline & Metric \\
\midrule
Sparse-data inverse PINN & sparse interpolation and wrong-parameter SIR & full S/I/R mean squared error \\
PIDL missing mechanism & known SIR with the missing term removed & full S/I/R mean squared error \\
Controlled malware mitigation & no control, fixed controls, direct PINN, PMP-informed PINN, and rollout-optimized neural control & rollout objective and cumulative compromised devices \\
\bottomrule
\end{longtable}

For control experiments, inspect both the training diagnostic and the rollout comparison. A control network can reduce its training loss while still needing independent ODE-rollout validation.

\section{Extension Rules}
\begin{enumerate}
\item If observations are noisy, add a validation split and report parameter uncertainty.
\item If the missing mechanism may depend on interventions, give the correction network access to \((t,x,u)\), not only \((t,x)\).
\item If a control should be deployable as feedback, use a control network \(\pi(t,x)\) and train over multiple initial conditions.
\item If the model is high-dimensional, start with degree-level compartments before moving to full node-level or graph-neural PINNs.
\item If dynamics have impulses or jumps, train separate intervals or include jump residuals explicitly.
\end{enumerate}

\section{Minimal Reporting Checklist}
A clear Note 2 experiment should report:
\begin{itemize}
\item model equations and true or assumed parameters;
\item observed variables, observation times, and noise level;
\item collocation sampling strategy;
\item neural objects and architecture sizes;
\item every loss term and its weight;
\item learned parameters or learned controls;
\item independent checks such as held-out data, ODE-solver validation, or comparison with FBSM/control baselines.
\end{itemize}

\end{document}
